\newcounter{goparagraphcounter}
\setcounter{goparagraphcounter}{1}

\section{Geschäftsordnung}

\subsection{Präambel}
Diese Geschäftsordnung gilt für die Mitglieder nach \parref{sa:par:mitgliedschaft} der Satzung des Narrenbund Ostfildern e.V. Sie regelt die interne Arbeitsweise und die Aufgabenverteilung innerhalb des Vereins. 

\subsection{§ \thegoparagraphcounter\ Erlass, Änderung, Aufhebung und Bekanntmachung} \label{go:par:erlass_aenderung_aufhebung}
\begin{enumerate}
    \item \label{go:abs:kwyvk7} Diese Geschäftsordnung kann jederzeit durch den Vereinsausschuss erarbeitet werden.

    \item \label{go:abs:1rro4u} Erlass, Änderungen oder Aufhebung der Geschäftsordnung bedürfen der Genehmigung durch die Mitgliederversammlung gemäß \parabsref{sa:par:ordnungen}{sa:abs:zyb87y} der Satzung.

    \item \label{go:abs:8z4qs6} Änderungen werden nach Genehmigung durch die Mitgliederversammlung wirksam und allen Mitgliedern bekannt gemacht.
\end{enumerate}

\refstepcounter{goparagraphcounter}
\subsection{§ \thegoparagraphcounter\ Interne Aufgaben- und Zuständigkeitsverteilung} \label{go:par:intern_aufgaben_zustaendigkeit}
\subsubsection{1. Vorstand}
Ergänzend zu den im \parabsref{sa:par:vorstand}{sa:abs:vln3it} der Satzung des Narrenbund Ostfildern e.V. geregelten Aufgaben übernimmt der 1. Vorstand folgende weitere Aufgaben:
\begin{itemize}
    \item Führt gemeinsam mit dem 1. stellv. Vorstand und dem 2. stellv. Vorstand die Geschäfte des Vereins und vertritt den Verein nach innen und außen (gesetzl. Verantwortung, Vertragsunterzeichnung u.Ä.).

    \item Beschließt gemeinsam mit dem 1. stellv. Vorstand und dem 2. stellv. Vorstand über die Verwaltung und die Ausgaben der Mittel im Rahmen des Haushaltsplans.

    \item Entscheidet gemeinsam mit dem 1. stellv Vorstand über Vereinsveranstaltungen, Versammlungen, etc.

    \item Plant und leitet die Mitgliederversammlung, Ausschusssitzungen und Vorstandssitzungen.

    \item Ist Ansprechpartner für alle Ämter.

    \item Besucht gemeinsam mit dem 1. stellv. Vorstand die Zunftmeisterempfänge anderer Zünfte.

    \item Erstellt den Haushaltsplan in Zusammenarbeit mit dem 1. stellv. Vorstand und dem 2. stellv. Vorstand.
\end{itemize}

\subsubsection{1. stellv. Vorstand}
Ergänzend zu den im \parabsref{sa:par:vorstand}{sa:abs:vln3it} der Satzung des Narrenbund Ostfildern e.V. geregelten Aufgaben übernimmt der 1. stellv. Vorstand folgende weitere Aufgaben:
\begin{itemize}
    \item Führt gemeinsam mit dem 1. Vorstand und dem 2. stellv. Vorstand die Geschäfte des Vereins und vertritt den Verein nach innen und außen (gesetzl. Verantwortung, Vertragsunterzeichnung u.Ä.).

    \item Beschließt gemeinsam mit dem 1. Vorstand und dem 2. stellv. Vorstand über die Verwaltung und die Ausgaben der Mittel im Rahmen des Haushaltsplans.

    \item Entscheidet gemeinsam mit dem 1. Vorstand über Vereinsveranstaltungen, Versammlungen, etc.

    \item Ist Ansprechpartner für alle Ämter.

    \item Unterstützt und berät den 1. Vorstand bei allen brauchtumsfördernden, wirtschaftlichen und sozialen Fragen.

    \item Besucht gemeinsam mit dem 1. Vorstand die Zunftmeisterempfänge anderer Zünfte.

    \item Leitet bei Verhinderung des 1. Vorstands die Mitgliederversammlung, Ausschusssitzungen und Vorstandssitzungen.

    \item Wickelt gemeinsam mit dem 2. stellv. Vorstand die Busplanung für die Narrenausfahrten ab.

    \item Plant die Veranstaltungen während der Kampagne; zur Durchführung kann er weitere Mitglieder oder Arbeitsgruppen einbinden und beauftragen.

    \item Erstellt den Haushaltsplan in Zusammenarbeit mit dem 1. Vorstand und dem 2. stellv. Vorstand.
\end{itemize}

\subsubsection{2. stellv. Vorstand}
Ergänzend zu den im \parabsref{sa:par:vorstand}{sa:abs:vln3it} der Satzung des Narrenbund Ostfildern e.V. geregelten Aufgaben übernimmt der 2. stellv. Vorstand folgende weitere Aufgaben:
\begin{itemize}
    \item Führt gemeinsam mit dem 1. Vorstand und dem 1. stellv. Vorstand die Geschäfte des Vereins und vertritt den Verein nach innen und außen (gesetzl. Verantwortung, Vertragsunterzeichnung u.Ä.).

    \item Beschließt gemeinsam mit dem 1. Vorstand und dem 1. Vorstand über die Verwaltung und die Ausgaben der Mittel im Rahmen des Haushaltsplans.

    \item Führt die Vereinskassen und Konten.

    \item Klärt in Absprache mit dem Steuerberater Steuerrechtliche Fragen.

    \item Wickelt gemeinsam mit dem 1. stellv. Vorstand die Busplanung für die Narrenausfahrten ab.

    \item Erstellt den Jahresabschluss des jeweiligen Geschäftsjahres.

    \item Erstellt die Kassenberichte nach Veranstaltungen und präsentiert diese dem Vereinsausschuss. 
    
    \item Erstellt den Haushaltsplan in Zusammenarbeit mit dem 1. Vorstand und dem 1. stellv. Vorstand.
\end{itemize}

\subsubsection{Schriftführer}
\begin{itemize}
    \item Führt die Protokolle aller Sitzungen und Versammlungen.

    \item Informiert die Mitglieder durch Rundschreiben und ähnliches über wichtige vereinsinterne Ereignisse bzw. Vorhaben.

    \item Verwaltung der Social-Media-Kanäle des Vereins (Instagram, TikTok) und regelmäßiges posten von aktuellen Inhalten und evtl. Werbung.

    \item Verwaltung der Internetseite.
    
    \item Betreuung und Organisation von Beiträgen des Vereins in der Stadtrundschau.
\end{itemize}

\subsubsection{Festwart}
\begin{itemize}
    \item Allg. Organisation der Veranstaltungen des Narrenbund Ostfildern e.V.

    \item Planung von Ausflügen o.Ä. Aktivitäten außerhalb der Kampagne.
    
    \item Erstellung von Arbeitsdienstlisten.

    \item Erstellung von Preislisten nach Abstimmung mit dem Ausschuss.

    \item Koordination Aufbau und Abbau sowie Dekoration.

    \item Koordination von Arbeitsdiensten.
\end{itemize}

\subsubsection{Häswart}
\begin{itemize}
    \item Plant und leitet die Häskontrolle.

    \item Verwaltung der Häsmaterialien im Lager.

    \item Materialbestellung in Abstimmung mit dem Ausschuss

    \item Überwachung und Kontrolle der Häsordnung.

    \item Kontakt zu Maskenschnitzer, Schneiderei u.Ä.
\end{itemize}

\subsubsection{Jugendwart}
\begin{itemize}
    \item Stellt gemeinsam mit der Vorstandschaft die Einhaltung des Jugendschutzgesetzes sicher.

    \item Leitung und Koordination der Jugendaktivitäten im Verein. 

    \item Förderung des Zusammenhalts innerhalb der Jugendgruppe.

    \item Pflege der Kommunikation mit Eltern und Jugendlichen.

    \item Förderung der Persönlichkeitsentwicklung und sozialen Kompetenzen der Jugendlichen.

    \item Sicherstellung der Einhaltung von Vereinswerten und Traditionen.
\end{itemize}

\subsubsection{Vereinsausschuss}
\begin{itemize}
    \item Berät den Vorstand in allen wesentlichen Vereinsangelegenheiten.
    
    \item Unterstützt die Organisation und Durchführung von Vereinsveranstaltungen.
    
    \item Überwachung der Satzung und der Ordnungen des Vereins.
    
    \item Schlichtung und Aussprache von Problemen innerhalb des Vereins
    
    \item Aufnahme und Ausschluss von Vereinsmitgliedern.
\end{itemize}

\refstepcounter{goparagraphcounter}
\subsection{§ \thegoparagraphcounter\ Wahlturnus} \label{go:par:wahlturnus}
\begin{enumerate}
    \item \label{go:abs:e1ey4m} Für die Vorstandsämter wird folgender Wahlturnus festgelegt:

    \vspace{0.5cm}
    \noindent
    \begin{tabular}{p{5cm}p{5cm}p{5cm}}
    \textbf{1. Jahr} &
    \textbf{2. Jahr} &
    \textbf{3. Jahr} \\

    \begin{itemize}
        \item 1. Vorstand
    \end{itemize} &
    \begin{itemize}
        \item 2. stellv. Vorstand
    \end{itemize} &
    \begin{itemize}
        \item 1. stellv. Vorstand
    \end{itemize} \\
    \end{tabular} \\

    \vspace{0.5cm}

    \item \label{go:abs:nvhuj8} Wird ein Mitglied des Vorstands gemäß \parabsref{sa:par:vorstand}{sa:abs:fiifa7} der Satzung neu gewählt, gilt dies nur für die verbleibende Amtszeit des ausgeschiedenen Vorstandsmitglieds, sodass der bestehende Wahlturnus des Vereins unberührt bleibt. \\

    \item \label{go:abs:fhc5rj} Für die Ausschussämter wird folgender Wahlturnus festgelegt:

    \vspace{0.5cm}

    \begin{tabular}{p{5cm}p{5cm}}
    \textbf{1. Jahr} &
    \textbf{2. Jahr} \\

    \begin{itemize}
        \item Häswart
        \item Jugendleiter
        \item Kassenprüfer 1
    \end{itemize} &
    \begin{itemize}
        \item Schriftführer
        \item Festwart
        \item stellv. Jugendleiter
        \item Kassenprüfer 2
    \end{itemize} \\
    \end{tabular}
\end{enumerate}

\refstepcounter{goparagraphcounter}
\subsection{§ \thegoparagraphcounter\ Ausschüsse} \label{go:par:ausschuesse}
\begin{enumerate}
    \item \label{go:abs:pdl45j} Der Vereinsausschuss kann zur Aufgabenerledigung einzelne oder mehrere Mitglieder beauftragen.

    \item \label{go:abs:fs67un} Die Berufung erfolgt nach Bedarf und ist nicht an Inhalte und Aufgabenstellungen gebunden. Der Vereinsausschuss entscheidet insoweit nach freiem Ermessen.

    \item \label{go:abs:7m0wd4} Die beauftragten Mitglieder haben keine Entscheidungsbefugnis. Sie dienen der Beratung und Meinungsbildung für den Vereinsausschuss und bereiten Entscheidungen vor.
\end{enumerate}

\refstepcounter{goparagraphcounter}
\subsection{§ \thegoparagraphcounter\ Mitgliedschaft} \label{go:par:mitgliedschaft}
\begin{enumerate}
    \item \label{go:abs:2o7f6t} Die Aufnahme neuer Mitglieder erfolgt ausschließlich als Probeläufer durch schriftliche oder mündliche Bestätigung des Vorstandes.

    \item \label{go:abs:e6pthy} Bei Eintritt erhält jeder neue Probeläufer die jeweils gültige Satzung sowie das vollständige Ordnungswerk des Vereins.

    \item \label{go:abs:h8muwn} Die Aufnahme von Personen unter 16 Jahren ist in der Regel nur möglich, wenn mindestens ein gesetzlicher Vertreter ebenfalls ordentliches Mitglied des Vereins ist bzw. wird. Der Ausschuss kann in begründeten Fällen Ausnahmen zulassen.

    \item \label{go:abs:daoe1b} Jeder aufgenommene Probeläufer absolviert ein Probejahr. Der Ausschuss kann in begründeten Fällen das Probejahr verlängern.

    \item \label{go:abs:cs55q5} Während des Probejahres erhält das Mitglied ausschließlich Vereinskleidung; ein Häs wird in dieser Zeit nicht angefertigt.

    \item \label{go:abs:4t9fyx} Probeläufer müssen sich aktiv an sämtlichen Vereinsaktivitäten während und außerhalb der Kampagne beteiligen.

    \item \label{go:abs:uvjuiz} Nach erfolgreichem Abschluss des Probejahres entscheidet der Ausschuss über die endgültige Aufnahme als Mitglied gemäß \parabsref{sa:par:erwerb_erloeschen_mitgliedschaft}{sa:abs:66yn4s} der Satzung. Erst nach der endgültigen Aufnahme als Mitglied ist die Anfertigung des jeweiligen Häs zulässig.

    \item \label{go:abs:pd9q04} Die Narrentaufe erfolgt, sobald das Mitglied ein vollständiges Häs besitzt, in der darauffolgenden Kampagne.

    \item \label{go:abs:qbkq7j} Der reguläre Zeitraum zur Aufnahme als Probeläufer liegt zwischen Aschermittwoch und dem 31. August. Ausnahmen hiervon können vom Ausschuss beschlossen werden.

    \item \label{go:abs:0fzfll} Bei Austritt aus dem Verein sind das Vereinswappen sowie die dem Mitglied zugewiesene Laufnummer unverzüglich, spätestens jedoch innerhalb von 14 Tagen ab dem Zeitpunkt der Austrittserklärung, unaufgefordert an den Häswart zurückzugeben. Davon unberührt bleiben die Bestimmungen zu Häs und Eigentum gemäß der Häsordnung.
\end{enumerate}

\refstepcounter{goparagraphcounter}
\subsection{§ \thegoparagraphcounter\ Stellung der einzelnen Maskengruppen und Abteilungen} \label{go:par:maskengruppen_abteilungen}
\begin{enumerate}
    \item \label{go:abs:jr0t2g} Dem Narrenbund Ostfildern e.V. gehören folgende Gurppen an:
    \begin{enumerate}
        \item Rotbartle
        \item Kropfschell
        \item Glengabachdeifl
    \end{enumerate}
    
    \item \label{go:abs:vwhkz5} Die Maskengruppen sind organisatorische Untergliederungen des Vereins. Die Verantwortung für Finanzen, Vereinsführung, Verwaltung, Öffentlichkeitsarbeit sowie für sämtliche übergeordneten Vereinsangelegenheiten liegt ausschließlich beim Verein bzw. dem Vereinsausschuss. Den Maskengruppen kommt keine eigenständige rechtliche oder finanzielle Entscheidungsbefugnis zu.

    \item \label{go:abs:uq4co2} Der Verein tritt nach außen als Einheit auf. Veranstaltungen, Umzüge und sonstige offizielle Termine werden grundsätzlich gemeinsam besucht. Einzelne Maskengruppen handeln hierbei nicht eigenständig und nehmen ohne Abstimmung mit dem Vorstand oder dem zuständigen Ausschuss nicht separat an Veranstaltungen teil.

    \item \label{go:abs:hhschu} Die Zahl der aktiven Hästräger einer Maskengruppe soll 30 Personen nicht überschreiten. Ist die maximale Gruppenstärke erreicht, wird für weitere Interessenten eine Warteliste geführt. Über die Reihenfolge, Aufnahme und mögliche Ausnahmen entscheidet der Ausschuss.
\end{enumerate}

\refstepcounter{goparagraphcounter}
\subsection{§ \thegoparagraphcounter\ Gruppenwechsel innerhalb des NBO} \label{go:par:gruppenwechsel}
\begin{enumerate}
    \item \label{go:abs:zwvpr7} Möchte ein Hästräger des Narrenbund Ostfildern e.V. zu einer anderen Maskengruppe wechseln ist wie folgt zu verfahren:
    \begin{enumerate}
        \item Ein Wechsel der Maskengruppe ist nur auf Antrag des Hästrägers möglich. Der Antrag ist schriftlich an den Vorstand zu richten und muss die gewünschte Zielgruppe enthalten.
        
        \item Der Vereinsausschuss entscheidet in der folgenden Ausschusssitzung über den Antrag mit einfacher Mehrheit.

        \item Wird dem Antrag zugestimmt, erfolgt der Wechsel mit Wirkung zum nächstmöglichen Zeitpunkt, der organisatorisch vertretbar ist.

        \item Sofern keine Ausnahmegenehmigung des Ausschusses vorliegt, ist ein Wechsel erst zulässig, wenn der Hästräger mindestens fünf Jahre der bisherigen Maskengruppe angehört hat.

        \item Der Ausschuss kann in begründeten Fällen Ausnahmen zulassen (z. B. gesundheitliche Gründe, familiäre Gründe, interne Umstrukturierungen).
    \end{enumerate}
\end{enumerate}

\refstepcounter{goparagraphcounter}
\subsection{§ \thegoparagraphcounter\ Sprunggeld} \label{go:par:sprunggeld}
\begin{enumerate}
    \item \label{go:abs:e1si6l} Zusätzlich zum regulären Mitgliedsbeitrag wird einmal jährlich vor Beginn der Kampagne ein gesonderter Betrag, das sogenannte \textit{Sprunggeld}, von allen aktiven Hästrägern erhoben.

    \item \label{go:abs:84buky} Das Sprunggeld dient der Finanzierung der während der Kampagne anfallenden Kosten, insbesondere für Busfahrten sowie benötigte Materialien.

    \item \label{go:abs:yaloam} Der Verein bezuschusst die Busfahrten, um die Teilnahme für die Mitglieder finanziell im zumutbaren Rahmen zu halten. Die Höhe des Zuschusses und damit die Höhe des jährlich zu entrichtenden Sprunggeldes werden vom Ausschuss festgelegt.

    \item \label{go:abs:cp3v21} Außerordentliche Mitglieder, die an einer Busfahrt teilnehmen möchten, haben den jeweils geltenden Einzelpreis für die Busfahrt vollständig selbst zu tragen.

    \item \label{go:abs:rflvyp} Eine Rückerstattung des Sprunggeldes ist in jedem Fall ausgeschlossen, unabhängig vom Grund der Nichtteilnahme an Veranstaltungen oder Busfahrten.
\end{enumerate}

\refstepcounter{goparagraphcounter}
\subsection{§ \thegoparagraphcounter\ Spenden} \label{go:par:spenden}
\begin{enumerate}
    \item \label{go:abs:fs63cc} Etwaige Spenden fließen dem Narrenbund Ostfildern e.V. zu und werden ordnungsgemäß in der Vereinsbuchhaltung erfasst. Der Verein stellt die entsprechenden Spendenbescheinigungen bzw. -quittungen innerhalb des Kalenderjahres aus, in dem die Spenden zugegangen sind.
\end{enumerate}

\refstepcounter{goparagraphcounter}
\subsection{§ \thegoparagraphcounter\ Verhaltensregeln} \label{go:par:verhaltensregeln}
\begin{enumerate}
    \item \label{go:abs:gqzfdq} Jeder Hästräger ab dem 14. Lebensjahr ist verpflichtet, im vollständigen Häs, entsprechend der jeweils gültigen Häsordnung zu erscheinen.

    \item \label{go:abs:egp5vc} Das Ausleihen sowie der Tausch der Maske bzw. des Häs unter den Mitgliedern ist nicht gestattet. Ausnahmen können vom Vorstand im Einzelfall genehmigt werden.

    \item \label{go:abs:rdu1j3} Der Konsum illegaler Drogen ist grundsätzlich untersagt.

    \item \label{go:abs:lqdpi4} Der Konsum von Cannabis in Vereinskleidung oder im Häs ist untersagt.

    \item \label{go:abs:j3n0vb} Der Alkoholkonsum ist bei der Teilnahme an Veranstaltungen auf ein erträgliches Maß zu beschränken. Für Jugendliche unter 18 Jahren gelten die Bestimmungen des Jugendschutzgesetzes. Bei einem Verstoß kann die Teilnahme an der Veranstaltung durch ein Vorstandsmitglied untersagt werden.

    \item \label{go:abs:prjq6f} Bei einer selbstverschuldeten Schlägerei im Häs, wird eine 10 wöchige Sperre ausgesprochen. Der Vereinsausschuss kann nach eigenem Ermessen weitere Maßnahmen erlassen.

    \item \label{go:abs:7bp32s} Ein Fernbleiben an Umzügen und Veranstaltungen muss rechtzeitig einem Vorstandsmitglied mitgeteilt werden.

    \item \label{go:abs:zmmaom} Für jeden Hästräger ist die Teilnahme an 75 \% der Umzüge und Veranstaltungen verpflichtend. Familien zählen als eine Person. Sollte diese Grenze nicht erreicht werden, kann der Vereinsausschuss nach eigenem Ermessen Maßnahmen erlassen. Ausnahmen hiervon können vom Ausschuss beschlossen werden.

    \item \label{go:abs:hnnpec} Alle Hästräger haben pünktlich zum jeweiligen Aufstellungsplatz bei Umzügen zu erscheinen. Ein späteres Einsteigen in den laufenden Umzug ist grundsätzlich nicht vorgesehen.

    \item \label{go:abs:lubbg9} Jeder Hästräger hat dafür Sorge zu tragen, dass kein Zuschauer zu Schaden kommt. Besondere Rücksicht erfordern z.B. Kinder, ältere und behinderte Menschen, Schwangere, usw. Im Schadensfall muss umgehend ein Vorstandsmitglied verständigt werden. Bei Sachbeschädigung sowie Körperverletzung haftet ausschließlich der Hästräger.

    \item \label{go:abs:qq8m3q} Bei außerfasnachtlichen Veranstaltungen des Vereins gelten die gleichen Verhaltensregeln wie während der Kampagne.

    \item \label{go:abs:bp91kr} Jedes Mitglied verpflichtet sich, das in der Häsordnung beschriebene Häs, sowie die Maske sorgfältig zu behandeln und instand zu halten.
\end{enumerate}

\refstepcounter{goparagraphcounter}
\subsection{§ \thegoparagraphcounter\ Schlussbestimmungen} \label{go:par:schlussbestimmungen}
Bei Zuwiderhandlungen gegen die vorgenannten Bestimmungen können folgende Massnahmen
getroffen werden:
\begin{enumerate}
    \item \label{go:abs:fomeo1} Verwarnungen: mündliche Verwarnung / Abmahnung durch Ausschuss

    \item \label{go:abs:ug83mo} Sperren: Der Hästräger wird für maximal 10 Veranstaltungen gesperrt. Er darf in dieser Zeit sein Häs nicht tragen. Über die endgültige Dauer der Sperre entscheidet der Vereinsausschuss.

    \item \label{go:abs:17lqc7} Ausschluss: Der Hästräger wird aus der Gruppe / dem Verein ausgeschlossen, d.h. er muss sein Häs abgeben bzw. darf es nicht mehr tragen. Über den Ausschluss entscheidet der Ausschuss.

    \item \label{go:abs:4czlee} Bei Verstoß gegen das Satzung- und Ordnungswerk nach Austritt oder Ausschluss aus dem Verein, behält sich der Narrenbund Ostfildern e.V. ausdrücklich gerichtliche Schritte vor. 
\end{enumerate}

\refstepcounter{goparagraphcounter}
\subsection{§ \thegoparagraphcounter\ Inkrafttreten} \label{go:par:inkrafttreten}
\begin{enumerate}
    \item \label{go:abs:hsdf7s} Diese Geschäftsordnung wurde in der Ausschusssitzung vom dd. MMMM YYYY beschlossen und durch die Mitgliederversammlung vom dd. MMMM YYYY genehmigt. \\
    Sie ersetzt alle vorherigen Versionen.
\end{enumerate}